% SGA 3, Expose VI_B, Loop 1 English reconstruction.
% Scope: Corollary 6.5.5 and its proof in full, including source
% editor's note 66.
% Authority: Polo--Gille Exp6B-13oct24.pdf, component-local pages
% 48--49 (combined-reader pages 389--390; re-edition headers
% 378--379).
% Definition 6.6 begins on local page 49 and is excluded.
% OCR and Reinhold English Markdown are comparison witnesses only.
% This unit contains no diagrams.

\medskip
\phantomsection\label{cor:sga3-VIB-6.5.5}%
\noindent\textbf{推论 6.5.5.---}
\emph{设 \(G\) 为在 \(S\) 上局部有限型
且拟分离的 \(S\)-群概形, \(H\) 为 \(S\) 上的有限呈示平滑群
概形, 其纤维连通,
\(i:H\to G\) 为 \(S\)-群的单态射. 则:}
\begin{enumerate}
\item[\textup{a)}]\phantomsection\label{concl:sga3-VIB-6.5.5-a}
  \emph{\(\underline{\operatorname{Centr}}_G(H)\) 和
  \(\underline{\operatorname{Norm}}_G(H)\) 可由
  \(G\) 的闭子概形表出, 且在 \(G\) 上有限呈示.}
\item[\textup{a\(^{\prime}\))}]
  \phantomsection\label{concl:sga3-VIB-6.5.5-aprime}
  \emph{同样, 对 \(S\)-群概形的每个有限呈示单态射 \(j:K\to G\),
  \(\underline{\operatorname{Transp}}_G(H,K)\) 可由
  \(G\) 的闭子概形表出, 且在
  \(G\) 上有限呈示.}
\item[\textup{b)}]\phantomsection\label{concl:sga3-VIB-6.5.5-b}
  \emph{若 \(G\) 拟紧, 则存在一个整数
  \(n\geq 0\), 使得若 \(H_n\) 表示 \(H\) 的单位截面的第 \(n\) 阶
  无穷小邻域, 则}
  \[
    \underline{\operatorname{Centr}}_G(H)
      =\underline{\operatorname{Centr}}_G(H_n),
    \qquad
    \underline{\operatorname{Norm}}_G(H)
      =\underline{\operatorname{Norm}}_G(H_n),
  \]
  \[
    \underline{\operatorname{Transp}}_G(H,K)
      =\underline{\operatorname{Transp}}_G(H_n,K)
      =\underline{\operatorname{Transp}}_G(H_n,K_n).
  \]
\end{enumerate}

\emph{证明.}%
\footnote{\label{note:sga3-VIB-66}编者注: 下述论证
作了扩充. 至此也完成了 XI,
\SGARefLink{sga3:xi:theorem:6:8}{6.8}--6.11 的插入; 在下面的 \SGARefLink{def:sga3-VIB-6.6}{6.6} 中,
我们回到 VI\textsubscript{B} 的原文.}
首先注意, 关于 \(G\) 的假设蕴涵
单态射 \(H\to G\) 有限呈示
(EGA~IV\(_1\), 1.2.4 及 1.4.3), 对角浸入
\[
  \Delta_{G/S}:G\longrightarrow G\times_S G
\]
也有限呈示 (同上, 1.4.3.1).
因此取 \(H=K\), 即可将
\(\underline{\operatorname{Norm}}_G(H)\) 的情形化归为
传送子的情形.

根据 \hyperref[ex:sga3-VIB-6.2.4-e]{6.2.4\textup{(e)}},
对底概形 \(G\) 上的群概形
\[
  X=H_G=H\times_S G
\]
及下述子群概形 \(Y\),
应用 \hyperref[cor:sga3-VIB-6.5.4]{6.5.4}.

在 \(\underline{\operatorname{Transp}}_G(H,K)\)
的情形, 取 \(Y\) 为 \(K_G\) 在 \(G\)-群态射
\[
  H\times_S G\longrightarrow G\times_S G,
  \qquad
  (h,g)\longmapsto(ghg^{-1},g).
\]
下的逆像. 在
\(\underline{\operatorname{Centr}}_G(H)\) 的情形,
取 \(Y\) 为 \((G\times_S G)_G\) 的对角子群
\[
  \Delta_{G/S}(G)_G
\]
在 \(G\)-群态射
\[
  H\times_S G\longrightarrow G\times_S G\times_S G,
  \qquad
  (h,g)\longmapsto(h,ghg^{-1},g).
\]
下的逆像.

% 翻译统计: 源行数 91; 输出行数 92 (含本统计行); 行内/行间数学片段 44; 翻译占位 0. 底本疑点另报主线程.
